\chapter{Introduction}
\label{ch:introduction}

% TODO: write last (introductions are easier to write once the rest is in place).
% Structure mirrors the host lab's thesis format:
%   1.1 Context     -> background and motivation
%   1.2 Goals       -> objectives, research questions, contributions
%   1.3 Structure   -> chapter-by-chapter roadmap

This chapter introduces the context, motivations, and objectives of the
thesis. Section~\ref{sec:intro-context} frames the background of
multi-party speech-based conversational systems and AI moderation.
Section~\ref{sec:intro-goals} states the goals of the thesis, the
research questions, and summarises the contributions.
Section~\ref{sec:intro-structure} provides a roadmap of the remaining
chapters.

\section{Context}
\label{sec:intro-context}
% TODO: background on multi-party voice systems, motivation for AI moderation,
% gap that this work addresses (link to Chapter 2).

\section{Goals of the Thesis}
\label{sec:intro-goals}
% TODO: state goals + research questions + summary of contributions.
% RQ1 — How can a speech-based conversational system be architected to support
%       multi-party group discussions with AI moderation?
% RQ2 — What strategies can the AI moderator adopt to intervene effectively
%       without compromising the dynamics of the group?
% RQ3 — How does the presence/absence of the AI moderator affect group
%       performance, participation balance, and user experience in a
%       collaborative survival task?

\section{Structure of the Thesis}
\label{sec:intro-structure}
% TODO: one short paragraph per chapter that follows.
